\section{Background, related results, and the main theorem}
\label{sec:introduction}

Throughout, rings are nonzero, commutative, and have an identity,
and dimension means Krull dimension. The notation $R[[x]]$ denotes
the ring of formal power series in one indeterminate over $R$.
For an ideal $I\subseteq R$, we write
\[
 I[[x]]=\left\{\sum_{j\ge0}a_jx^j:a_j\in I\right\}.
\]
This is the kernel of the ring homomorphism $R[[x]]\longrightarrow(R/I)[[x]]$; it need not
equal the extended ideal $IR[[x]]$. We write $\N=\{1,2,\ldots\}$
and $\N_0=\N\cup\{0\}$.

The dimension theory of formal power series differs substantially
from that of polynomial rings. For a finite-dimensional ring $R$,
Seidenberg's bounds \cite{Seidenberg1953} give
\[
 \dim R+1\le\dim R[x]\le2\dim R+1.
\]
The lower bound also holds for $R[[x]]$, and equality holds when $R$
is Noetherian. In general, however, $R[[x]]$ can have infinite
dimension, and even its finite dimension need not satisfy the
polynomial upper bound.

An ideal $I$ is of \emph{strong finite type} (SFT) if there exist a
finitely generated ideal $J\subseteq I$ and an integer $e\ge1$ such
that $a^e\in J$ for every $a\in I$. A ring is SFT if each of its
ideals is SFT. Arnold \cite{Arnold1973a} proved that
\begin{equation}\label{eq:arnold}
 \dim R[[x]]<\infty\quad\Longrightarrow\quad R\text{ is SFT};
\end{equation}
see also Roitman~\cite{Roitman2015}. Condo, Coykendall, and
Dobbs~\cite[Corollary~2.2]{CondoCoykendallDobbs1996} proved that,
for a zero-dimensional SFT ring $R$,
\begin{equation}\label{eq:zero-classical}
 \dim R[[x_1,\ldots,x_s]]=s\qquad(s\ge1).
\end{equation}
These results account for the exceptional finite pair $(0,1)$.

Several important classes have more restrictive dimension formulas.
For a finite-dimensional Pr\"ufer domain $D$, Arnold's theorem
gives $\dim D[[x]]=\dim D+1$ precisely when $D$ is SFT, and
infinite dimension otherwise~\cite[Theorem~3.8]{Arnold1973b}.
More generally, if $D$ is an SFT Pr\"ufer domain of positive finite
dimension $d$, then (see \cite[Theorem~3.6]{Arnold1982})
\[
 \dim D[[x_1,\ldots,x_s]]=sd+1.
\]
The SFT hypothesis alone does not ensure finite power-series
dimension: Coykendall~\cite{Coykendall2002} constructed a
one-dimensional SFT domain whose power series ring has infinite
dimension.

Arnold~\cite{Arnold1981} had already shown that a finite
power-series dimension need not equal the coefficient dimension
plus one. Kang and Park~\cite{KangPark2009} constructed finite-dimensional
counterexamples to the bound $\dim R[[x]]\le2\dim R+1$.
Their mixed-extension formula yields, for suitable SFT Pr\"ufer
domains $D$ of dimension $d$ and $R=D[z_1,\ldots,z_p]$,
\[
 (\dim R,\dim R[[x]])=(d+p,\,d(p+1)+1).
\]
For example, $d=4$ and $p=3$ give $(7,17)$.
Thus finite violations of the polynomial bound are established
prior results. The construction below instead keeps the coefficient
dimension equal to one and prescribes every finite power-series
dimension at least two.

Our method also relates to the study of subrings of power series
rings defined by bounds on the degrees of their coefficients; see
Kang and Toan~\cite{KangToan2015a,KangToan2015b}.
Here the coefficients are rational functions, and the filtration
controls both their denominators and their numerator degrees.
We prove the required division estimate and Noetherianity directly.

\begin{theorem}\label{thm:classification}
Let $n,m\in\N_0$. There exists a nonzero commutative ring $R$ with
identity such that
\[
 \dim R=n\qquad\text{and}\qquad\dim R[[x]]=m
\]
if and only if
\[
 (n,m)=(0,1)\qquad\text{or}\qquad 1\le n<m.
\]
\end{theorem}

The essential realization statement is stronger in dimension one.

\begin{theorem}\label{thm:local}
For every integer $m\ge2$, there exists a one-dimensional local
domain $T$ such that $\dim T[[x]]=m$.
\end{theorem}

For $m=2$ one may take $T=k[[c]]$, where $k$ is a field.
For $m=r+2\ge3$, Sections~\ref{sec:filtration}--\ref{sec:dimension}
construct a local domain $T_r$ and prove
\begin{equation}\label{eq:target}
 \dim T_r=1,\qquad \dim T_r[[x]]=r+2.
\end{equation}
Section~\ref{sec:classification} deduces
\cref{thm:classification}. Its higher-dimensional realizations use
finite direct products; the theorem is stated for rings without a
domain hypothesis.
